\chapter{Fine-tuning}

Fine-tuning процесот е организиран така што подоцна може да се додадат и споредат повеќе модели.
Податоците се делат на training, validation и test дел, а test примерите се чуваат за конечна споредба.

Тековниот експеримент користи LoRA, метод за параметарски ефикасно дообучување со low-rank матрици \cite{hu2021lora}.
QLoRA, кој го комбинира овој пристап со квантизиран основен модел, е планиран како следна варијанта за споредба \cite{dettmers2023qlora}.

\begin{table}[ht]
  \centering
  \small
  \begin{tabular}{llll}
    \toprule
    Модел & Метод & Епохи & Статус \\
    \midrule
    Qwen3-0.6B & LoRA, rank 16 & 3 & Обучен адаптер \\
    Друг модел & Не е утврден & Не е утврден & Планиран \\
    \bottomrule
  \end{tabular}
  \caption{Fine-tuning експерименти}
\end{table}

За Qwen3-0.6B се користени batch size 2, gradient accumulation 16, maximum sequence length 512 и learning rate $2\times10^{-4}$.
Обуката е извршена со bfloat16 precision на NVIDIA GeForce RTX 5080.
Зачувани се LoRA адаптерот, tokenizer-от и конфигурацијата потребна за негово повторно вчитување.
